%% ---------------------------------------------------------------------------
%% COMSNETS submission draft.
%%
%% Source material: design-docs/scheduler-study.md (primary),
%% design-docs/scheduler-design.md and design-docs/simulator-design.md
%% (architecture and simulator detail). Every number in Section VI is
%% reproducible from the scripts named in Table REPRO.
%%
%% Prepared for DOUBLE-BLIND review: the author block is anonymised and the
%% artifact section names no repository. Restore both for camera-ready.
%% Page budget is 8; see paper/README.md before adding to Section VII.
%% ---------------------------------------------------------------------------
\documentclass[conference]{IEEEtran}
\IEEEoverridecommandlockouts

\usepackage{cite}
\usepackage{amsmath,amssymb}
\usepackage{booktabs}
\usepackage{multirow}
\usepackage{graphicx}
\usepackage{url}

% Stop "OpenAirInterface" breaking as "Ope-nAirInterface".
\hyphenation{Open-Air-Inter-face}

% Tables in this paper are dense; give them a little air.
\renewcommand{\arraystretch}{1.15}

\begin{document}

\title{An Engineering Study of a Two-Tier QoS-Aware Scheduler\\
for Single-Cell Private 5G}

% Double-blind submission: anonymised. Restore the real author block and
% affiliations for camera-ready.
\author{\IEEEauthorblockN{Anonymous Author(s)}
\IEEEauthorblockA{Paper submitted for double-blind review}}

\maketitle

\begin{abstract}
A private 5G cell on a factory floor carries traffic whose classes want
qualitatively different things: guaranteed-bitrate uplink video, motion
control with millisecond deadlines, and elastic best-effort. Proportional
fair (PF), the de facto default and the scheduler OpenAirInterface ships,
represents neither a rate contract nor a deadline. Whether to replace it is
an engineering decision, and this paper reports one made carefully. We build
a two-tier scheduler---a network-utility-maximisation program setting
per-flow rate targets on a ${\sim}1$\,s horizon, and a per-slot
drift-plus-penalty tracker realising them under a symbol-accurate PRB and
PDCCH budget---and evaluate it against RR, PF and a class-aware gradient
baseline.

We make no strong novelty claim for the design; it composes known parts, and
we describe it because the composition is instructive. Our contribution is
what building and measuring it carefully produced. First, an independent
reproduction, in a simulator sharing no code with prior work, of two results
the field has reported separately: configured grants dominate when the
control channel binds (30/30 deadlines met against PF's 2/30), and
deadline-blind scheduling fails silently (8/8 against 5/8, with PF's misses
invisible in mean throughput). Second, and we think more useful: three
modelling shortcuts that \emph{manufactured or destroyed} results in our own
evaluation, each quantified against its corrected version. In the worst
case a headline finding---that our scheduler protected multi-flow uplink UEs
where the baselines collapsed them---disappeared entirely once the uplink
transport block was filled by the UE rather than the scheduler, as 3GPP
requires. Third, an adoption decision with its reasoning: which regimes the
design wins, which metric it loses on and why that is a metric choice rather
than a failure, and an exact feasibility test the scheduler already computes
that makes a graceful fallback possible without regime guessing. Simulator
and scripts are released.
\end{abstract}

\begin{IEEEkeywords}
private 5G, MAC scheduling, quality of service, network utility
maximisation, Lyapunov optimisation, configured grants, OpenAirInterface
\end{IEEEkeywords}

% ===========================================================================
\section{Introduction}
% ===========================================================================

A private 5G network on a factory floor carries a workload that public
mobile broadband does not. Mobile robots stream machine-vision and LIDAR
uplink at a sustained rate; motion-control loops exchange small payloads
under packet delay budgets (PDB) of a few milliseconds; firmware pushes and
log uploads soak up whatever is left. These classes do not simply want
\emph{more} throughput---they want \emph{different things}. A best-effort
flow wants its share of the residue. A guaranteed-bitrate (GBR) flow wants
a specific rate or nothing of value. A delay-critical flow wants its bytes
before a deadline, after which they are worthless.

The scheduler that ships in OpenAirInterface (OAI), and the de facto
default across LTE and NR, is proportional fair (PF)~\cite{jalali2000}.
PF ranks users by instantaneous rate over smoothed average rate; its fixed
point maximises $\sum_i \log R_i$~\cite{kelly1998} but has no representation of a
rate contract or a deadline.

A substantial literature adds QoS-awareness by modifying the PF metric with
an urgency term---M-LWDF~\cite{andrews2001}, EXP/PF and the
Exp-rule~\cite{shakkottai2002}, and the variants surveyed
in~\cite{capozzi2013}. A parallel thread casts allocation as network
utility maximisation (NUM)~\cite{kelly1998} or as Lyapunov drift-plus-penalty
control~\cite{neely2010,tassiulas1992,stolyar2005}. The question we ask is when is 
it worth for an operator to use a more sophisticated two level scheduler which
might involve an LP or RL type solver, a control loop, and a
non-trivial port into a real-time MAC. \emph{When is that worth paying
for?}

This paper answers that question for a private-5G factory target, where our two level
scheduler uses a Linear Programming Formulation to determine rate allocations
at a coarse time granularity and the MAC scheduler uses a simpler delay plus drift
based formulation to meet those rate targets.

\subsection*{Contributions}

We state the shape of this paper plainly, because it is not the usual one.
The scheduler composes established parts---network utility
maximisation~\cite{kelly1998} over a slow horizon feeding Lyapunov
drift-plus-penalty~\cite{neely2010} per slot, the pattern behind
NVS~\cite{kokku2012} and the O-RAN RIC split~\cite{oran}. We claim no
novelty for it, and present it (Section~\ref{sec:design}) because the
composition, and the places it goes wrong, are instructive. What we offer
instead:

\begin{enumerate}
\item \textbf{Independent reproduction of two results the field reports
separately} (Section~\ref{sec:results}), in a simulator sharing no code or
abstractions with the work that first reported them. Configured grants
dominate once the control channel binds---30/30 deadlines met against PF's
2/30---which Larra\~{n}aga et al.~\cite{larranaga2023cg} report from
ns-3/5G-LENA. And deadline-blind scheduling fails \emph{silently}: PF meets
5/8 deadlines while posting an 86\% mean that reads healthy on a dashboard.
Reproduction is not novelty, but in a field where most schedulers are
evaluated once, in one simulator, by their authors, it is worth having.

\item \textbf{Three modelling shortcuts that manufactured or destroyed
results in our own evaluation} (Section~\ref{sec:fidelity}), each reported
with the corrected measurement beside it. The sharpest: our scheduler
appeared to protect multi-flow uplink UEs that the QoS-blind baselines
collapsed to ${\sim}3\%$ of contract---a clean demonstration of QoS-aware
multiplexing, until the uplink transport block was filled by the \emph{UE}
rather than the scheduler, as TS 38.321 requires. The baselines then reach
100/90/92\% unaided; the protection was the UE's prioritised-bit-rate
configuration all along. We report all three because the shortcuts are
natural ones, the errors were ours, and the corrections changed conclusions.

\item \textbf{An adoption decision with its reasoning}
(Section~\ref{sec:adopt}). The design wins unconditionally on two mechanisms
PF structurally lacks, and loses on a knife-edge contract-count metric under
GBR infeasibility---because a threshold metric rewards concentrating a
shortfall and this design spreads it. We show that no knob recovers the
difference, that a scheduler-level fallback to PF would surrender the
mechanisms that work, and that the strategic tier already computes an exact
feasibility test making graceful degradation possible without regime
guessing.
\end{enumerate}

We also report two negative results---an adaptive dual-ascent GBR penalty and
a spectral-efficiency tilt, both implemented, measured and rejected---and
release the simulator, the scheduler library and one script per result table.

% ===========================================================================
\section{Background and Related Work}
% ===========================================================================

\subsection{Channel-opportunistic scheduling}
Round Robin is channel-blind and wastes capacity on faded users; Max-C/I
maximises cell throughput and starves the cell edge; PF~\cite{jalali2000}
interpolates via $r_{\mathrm{inst}}/R_{\mathrm{avg}}$ and is the practical
default. All three share one limitation: none can be told ``this flow needs
8\,Mbps'' or ``this packet dies in 10\,ms.''

\subsection{QoS-aware single-tier schedulers}
M-LWDF~\cite{andrews2001} multiplies the PF metric by head-of-line delay
and a weight; the Exp-rule~\cite{shakkottai2002} uses an exponential
urgency term. These handle \emph{delay} well. For \emph{rate} contracts
they share a subtler defect: a multiplicative urgency metric of the form
$(r_{\mathrm{inst}}/R_{\mathrm{avg}})\cdot(1+w\cdot\mathrm{deficit})$ has
an equilibrium that equalises a weighted $1/R_{\mathrm{avg}}$ across flows.
It cannot drive delivered rate to an \emph{arbitrary, unequal} target
vector however large $w$ is. This is not only an analytical objection.
Monghal et al.~\cite{monghal2008qos} weight their frequency-domain metric
by $\max(1, \mathrm{TBR}/R_n)$ for users below a target bit rate and
report that the weighting ``is not efficient enough to allow 95\% of the
users to reach their TBR''---the saturation, measured, in the setting it
was designed for. This saturation argument is why our Tier-2 is
a drift-plus-penalty integrator rather than a tuned PF variant; we include
a class-aware ``Gradient'' baseline of exactly this shape to make the
difference measurable.

\subsection{Optimisation-based and hierarchical RRM}
NUM~\cite{kelly1998} states allocation declaratively as
$\max \sum_i U_i(r_i)$ under capacity constraints, but yields a
\emph{rate vector}, whereas a wireless scheduler must decide \emph{which UE
to serve this slot} over a stochastic channel. Drift-plus-penalty
control~\cite{neely2010} and greedy primal-dual
scheduling~\cite{stolyar2005} address the per-slot problem but need a
target to track. The natural composition---a slow declarative program
feeding a fast tracker---is the pattern behind NVS~\cite{kokku2012}, the
O-RAN non-RT/near-RT RIC split~\cite{oran}, and two-timescale IIoT resource
allocation that likewise places Lyapunov control on the slow scale and a
fast solver beneath it~\cite{he2022twotimescale}. \textbf{The pattern is
not our contribution}; its concrete instantiation against a symbol-accurate
NR resource model, and the empirical characterisation of when it pays, are.

\subsection{Learning-based schedulers}
Most recent work in this space is learned rather than declarative. xSlice
places an online deep-RL agent on the near-RT RIC and adapts MAC-layer
allocation against a combined throughput/latency/reliability
objective~\cite{yan2026xslice}, evaluated---unlike this paper---on a real
O-RAN testbed. We deliberately do not take that route, and it is worth saying why. Kela et
al.~\cite{kela2024deep}, arguing from inside that literature, find that
deployable deep schedulers do not yet exist, citing real-time execution
cost, 3GPP compliance and poor generalisation across bandwidth, MIMO and
traffic configurations. For a deployment that must be commissioned and
certified by plant engineers, a convex program with an auditable objective
and a single fairness dial is a defensible alternative---and its failure
modes, as Section~\ref{sec:knapsack} shows, can be characterised
analytically rather than only measured.

\subsection{Industrial and private 5G}
Closest to our target is FLEX~\cite{kleinberger2026flex}, which is also
QoS-aware, also industrial, and also simulator-based, but varies the
\emph{TDD pattern}---adjusting the UL/DL split in flexible slots---where we
hold the pattern static and allocate within a fixed per-direction budget.
The two are complementary. On configured grants, Larra\~{n}aga et
al.~\cite{larranaga2023cg} contribute an independent open-source CG
implementation for ns-3/5G-LENA and an Industry 4.0 scheduling study; their
finding that CG policy dominates latency for time-critical industrial
traffic corroborates our Section~\ref{sec:pdcch} from a different
simulator. Our contribution there is accordingly not that CG helps, but
\emph{why} it helps twice over---bypassing the PDCCH budget and the BSR
round-trip simultaneously---and the self-gating rule that keeps it from
hurting when reservations cannot be honoured.

\subsection{The OAI baseline}
OAI's~\cite{oai} NR MAC schedules with a PF metric over per-UE grants. Our
\texttt{ProportionalFair} baseline is modelled on it, and all schedulers we
compare grant \emph{per UE} with one DCI each, so the comparison isolates
QoS-awareness rather than grant granularity. OAI is also the intended integration target; open private-5G O-RAN testbeds
built on it are now well established~\cite{villa2024x5g}.

% ===========================================================================
\section{System Model}\label{sec:model}
% ===========================================================================

We consider a single NR cell with a static TDD pattern. Time is slotted;
each slot offers a set of physical resource blocks (PRBs) in one direction,
with the special slot contributing partial DL and UL symbol counts. Let
$d(i)\in\{\mathrm{DL},\mathrm{UL}\}$ be flow $i$'s direction and $C_d$ the
capacity of direction $d$ in \emph{PRB-symbols per second} over the TDD
cycle. Per-UE SNR evolves as an AR(1) process in dB and maps through a
link-adaptation staircase to $\mathrm{SE}_i$, bits per
PRB-symbol, discounted by target BLER.

Each flow is uni-directional and carries a 5G QoS profile~\cite{ts23501}:
a class
$c(i)\in\{\mathrm{GBR},\mathrm{Delay},\mathrm{PF}\}$, a guaranteed flow
bit rate $G_i$ (GBR only), a packet delay budget $\mathrm{PDB}_i$, a
5QI-style \texttt{priority\_level}, and an optional slice identifier.
$D_i$ denotes offered demand.

Two control-plane budgets are modelled explicitly because both bind in
practice. First, a per-slot \textbf{PDCCH/CCE budget}: every dynamic grant
costs one DCI at an aggregation level that grows as SNR falls. Second, the
\textbf{uplink buffer-status-report round-trip}: a real gNB learns of UL
data only via a delayed BSR, so the scheduler sees a lagged view of UL
backlog. Semi-persistent scheduling (SPS, ``Configured
Grants''~\cite{ts38321}) bypasses both---a CG user transmits on a standing allocation with no DCI and no BSR.
Section~\ref{sec:pdcch} shows this double bypass is decisive.

% ===========================================================================
\section{Two-Tier Scheduler Design}\label{sec:design}
% ===========================================================================

\subsection{Tier-1: the strategic solve}

Tier-1 answers, every ${\sim}1$\,s, \emph{what rate should each flow
target?} Decision variables are target rates $r_i \ge 0$ (bps) with
shortfall slacks $s_i, s_\sigma \ge 0$ for GBR and slice floors. All solves
share one feasible set, parameterised by a per-flow floor:
\begin{align}
\mathcal{F}(\mathrm{floor}):\quad
&\textstyle\sum_{i: d(i)=d} r_i/\mathrm{SE}_i \le C_d && \forall d \label{eq:cap}\\
&r_i \le D_i && \forall i \label{eq:dem}\\
&r_i + s_i \ge G_i && \forall i{:}\, c(i){=}\mathrm{GBR} \label{eq:gbr}\\
&r_i \ge \mathrm{floor}_i && \forall i{:}\, c(i){=}\mathrm{GBR} \label{eq:hard}\\
&\textstyle\sum_{i\in\sigma} r_i/\mathrm{SE}_i + s_\sigma \ge \phi_\sigma C_d
  && \forall \sigma \label{eq:slice}
\end{align}
Constraint~\eqref{eq:cap} is where symbol-accurate accounting enters:
$r_i/\mathrm{SE}_i$ is the resource a target rate consumes, and $C_d$
counts special-slot symbols. Slice floors~\eqref{eq:slice} are capped at
the slice's own demand and kept soft, so an idle slice holds nothing and a
busy slice borrows freely---the allocation remains work-conserving.

\subsection{Shortfall before utility: an ordering, not a weight}
\label{sec:lex}

Ordering contracts ahead of utility is not new: Monghal et
al.~\cite{monghal2008qos} give below-target users absolute priority over
the rest, and Zaki et al.~\cite{zaki2011multi} give GBR bearers strict
priority over non-GBR. Both express it as a \emph{priority partition} over
a metric scheduler, which says who is served first but not how a shortfall
is divided among them. An optimisation formulation must answer that second
question too---and Section~\ref{sec:knapsack} is about the answer it gives.

The conventional formulation is a single objective,
\begin{equation}
\max \;\; \textstyle\sum_i w_{c(i)} \log(r_i + \epsilon) \;-\; \sum_i p_i s_i ,
\label{eq:single}
\end{equation}
with $p_i$ ``large'' so that GBR floors are effectively hard whenever
feasible. On our factory workload we measure the two terms at $718$ and
$2.4\times10^{10}$ respectively: a ratio of $3.3\times10^{7}$. That is not
a weighting. It is a \emph{lexicographic ordering expressed as a
magnitude}, and it has two consequences.

The first is numerical. Posed as~\eqref{eq:single} the program is badly
conditioned: on a three-flow instance small enough to solve in closed form,
two interior-point solvers both returned \texttt{optimal\_inaccurate},
disagreed by $3.6\times$ on one flow's rate, and under-served the
highest-weighted class by 28\% against the analytic optimum. Rescaling does
not fix it---a ratio between objective terms is a property of the model, not
the units. Stating the order directly does. Over $\mathcal{F}$ we solve
\begin{align}
\text{B1:}\quad S^\star &= \min \;\; \textstyle\sum_i p_i s_i + p_\sigma \sum_\sigma \mathrm{SE}_\sigma s_\sigma
  \label{eq:b1}\\
\text{B2:}\quad &\max \;\; \textstyle\sum_i w_{c(i)}\log(r_i+\epsilon)
  \;\;\text{s.t.}\;\; \eqref{eq:b1} \le S^\star
  \label{eq:b2}
\end{align}
after which both solvers agree to within 100\,bps of the analytic optimum.
A useful side effect: $p_i$'s absolute magnitude becomes irrelevant and
only the ratios across flows matter. Both slacks are converted to bits per
second before being weighed---a slice slack is natively in PRB-symbols---so
that $p_i$ and $p_\sigma$ are quoted in one currency. Compared raw, the
crossover between them sits at $p_i \cdot \mathrm{SE}_i$, which would make
slice-versus-GBR priority a function of the radio channel rather than of
policy.

\subsection{Soft GBR floors are a knapsack}\label{sec:knapsack}

The second consequence is structural. Once the penalty dominates,
\eqref{eq:b1} is $\min \sum_i (G_i - r_i)$ subject to
$\sum_i r_i/\mathrm{SE}_i \le C$: a \textbf{fractional knapsack}, whose
optimum is greedy in spectral efficiency and therefore a \emph{vertex} of
the feasible polytope. Flows are served in full or abandoned outright.
The solved targets are exactly that staircase: 100\% for the two highest-SE
flows and the four next, 50--67\% for a single boundary tier, and
\textbf{0\%} for the lowest-SE flow. Note it orders by SE, not SNR---ue2
(18\,dB) and ue4 (16\,dB) share an MCS step and receive identical targets.

Two controls confirm the mechanism. Sweeping $p$ over six decades leaves the
targets \emph{identical}---the solution is fixed by the knapsack structure,
not the penalty magnitude, so no choice of $p$ is a tuning opportunity. And
dropping the shortfall term entirely, the worst-SE flow receives 23\% of its
contract rather than $0$: the utility alone is SE-favouring but does not
abandon, and it is the shortfall term that converts a merely SE-favouring
allocation into a vertex one.

\textbf{This is a property of our formulation, not of QoS-aware
scheduling}, and we checked: none of the canonical GBR schedulers uses a
shortfall penalty. The literature encodes the contract either in a per-flow
metric under a priority
partition~\cite{monghal2008qos,zaki2011multi,ameigeiras2016qos} or in the
shape of the utility, kept \emph{concave} precisely so the allocation
problem has a unique solution~\cite{gora2014qos}---which cannot exhibit
this. We keep the penalty form because the lexicographic pair of
Section~\ref{sec:lex} expresses ``contracts first'' exactly where a concave
utility only approximates it, but a reader adopting it should know what it
carries.

\subsection{A max-min GBR floor}\label{sec:maxmin}

We therefore prepend a max-min satisfaction stage. Let
$\hat{G}_i = \min(G_i, D_i)$ be flow $i$'s \emph{reachable} contract; the
demand cap matters, or an under-offered GBR flow pins the level at its own
unreachable ratio and drags the set down. Stage~A solves
\begin{equation}
\begin{aligned}
t^\star = \max\, t \quad \text{s.t.} \quad & r\in\mathcal{F}(0), \;\; t\in[0,1],\\
& r_i \ge t\hat{G}_i \;\; \forall i{:}\, c(i){=}\mathrm{GBR}
\end{aligned}
\label{eq:maxmin}
\end{equation}
and sets $\mathrm{floor}_i = \alpha\, t^\star \hat{G}_i$ in~\eqref{eq:hard}.
$t^\star$ is the largest fraction of contract that \emph{every} GBR flow
can hold simultaneously; $t^\star = 1$ means the GBR set is jointly
feasible.

This is enabled by default, and the reason it is safe is that it
\textbf{self-disables}: whenever $t^\star = 1$ the floor binds nothing and
the result is identical to omitting the stage. It costs something only in
genuine GBR overload. $\alpha$ trades the guarantee against throughput,
with $\alpha = 0$ recovering the single-stage behaviour exactly.

All programs are posed in normalised units---rates as multiples of the
largest contract, capacity usage as a fraction of each direction's
budget---so every coefficient is $O(1)$. This is not cosmetic: posed in raw
bps, a unit-scale $t$ against $10^7$-scale rates made the solver report
\texttt{optimal} on a $t^\star$ that was \emph{non-monotone in capacity},
which is impossible for a max-min level.

\subsection{Tier-2: drift-plus-penalty rate tracking}

Each flow carries a virtual queue $Q_i$, a control accumulator in bits
rather than a buffer of data. Per slot of duration $\Delta$: grow
$Q_i \mathbin{+}= r_i\Delta$; clamp $Q_i \leftarrow \min(Q_i, \mathrm{ceil}_i)$;
serve; drain $Q_i \leftarrow \max(0, Q_i - \mathrm{delivered}_i)$. Because
$Q_i$ is a non-saturating integrator it can enforce an \emph{arbitrary}
target vector, which is precisely what the multiplicative metrics of
Section~II cannot do.

The clamp is subtle. $Q_i$ is bounded by what the flow legitimately should
have delivered over the trailing Tier-1 window $W$ but did not:
$\mathrm{ceil}_i = \max(0, \min(r_i W, A_i^W) - D_i^W)$ with $A^W, D^W$
bits arrived and delivered in the window. Using a \emph{windowed arrival
count} rather than instantaneous backlog is essential: a bursty video
flow's buffer momentarily empties between frames, and clamping to that
near-zero backlog erases its legitimate rate-tracking debt, letting
continuous flows starve bursty GBR ones. This was a real regression in our
implementation.

\subsection{Per-UE grants and the MAC multiplexer}

The 5G MAC grants per UE---one DCI, one transport block---so Tier-2 does
too. For Delay flows a head-of-line urgency bonus enters the queue before
ranking,
$\tilde{Q}_i = Q_i + w_\mathrm{d}(\mathrm{HoL}_i/\mathrm{PDB}_i)^{\kappa}\max_j Q_j$,
scaled by the \emph{system-wide} maximum queue rather than the flow's own
small target---otherwise a low-rate control flow could never out-compete
bulk. UEs are then ranked by
$M_u = \big(\sum_{i\in u} \tilde{Q}_i\big)\cdot \mathrm{SE}_u$, which is
simultaneously opportunistic and rate-tracking, and served greedily subject
to PRB and CCE budgets. A granted UE's transport block is filled across its
flows by logical-channel prioritisation~\cite{ts38321}: by
\texttt{priority\_level}, then by $\tilde{Q}_i$. The baselines use the same multiplexer (with plain
backlog as the tiebreak, since they carry no virtual queues), so all
schedulers are strictly comparable on DCI cost.

\subsection{Configured grants}
Periodic flows receive standing allocations sized to their contracted
floor, granted in priority tiers with proportional scale-back so that no
flow is dropped for being late in an enumeration order. A \emph{viability
floor} governs the mechanism: if a tier's reservations would be scaled
below a threshold, the tier runs dynamically instead---unless dropping it
would overrun the CCE budget, in which case SPS's zero-DCI property makes
it the lesser evil. SPS thus self-gates, engaging where it helps and
standing aside where fixed reservations would starve the dynamic pool.

% ===========================================================================
\section{Evaluation Methodology}\label{sec:method}
% ===========================================================================

\subsection{Simulator and fidelity discipline}
We evaluate in a purpose-built discrete-event simulator whose job is
\emph{comparative} questions, not absolute ones. The design principle is to
model, to good accuracy, every resource the scheduler competes for, and
nothing else. Modelled: the slot-by-slot PRB grid including the TDD special
slot; the per-slot PDCCH/CCE budget with variable DCI aggregation level;
the UL BSR round-trip as a per-flow delay plus Bernoulli loss, bypassed by
configured grants; CQI report delay and loss, with BLER computed from the
mismatch between the MCS the scheduler picked and the true SNR at
transmission; per-UE AR(1) SNR through an MCS staircase; per-flow buffers
with head-of-line timestamps and PDB expiry. Not modelled: LDPC and I/Q,
HARQ as a full state machine (a fixed delay plus BLER discount instead),
per-packet RLC segmentation, MU-MIMO, mobility and inter-cell interference.

The discipline cuts both ways, and the PDCCH budget is the case in point:
modelling it is what reveals configured grants to be decisive
(Section~\ref{sec:pdcch}), and omitting it would have led to the opposite
conclusion. The same held for the BSR round-trip, which was our largest
fidelity gap until we closed it; omitting it \emph{understates} SPS.

\subsection{Scenarios}
Three workloads isolate three bottlenecks. \texttt{factory\_robots}: 10
mobile robots, 24 flows---uplink camera and LIDAR (GBR), downlink
motion-control (Delay), best-effort uploads and a TCP flow---with an SNR
spread of 14--24\,dB and uplink demand roughly twice carrier capacity as
shipped. Carrier: 40\,MHz, numerology $\mu=2$, DSUUU TDD (uplink-heavy).
\texttt{sensor\_dense}: 30 small periodic uplink sensors with a 15\,ms
PDB, sized so the control channel binds before the data channel.
Carrier: 30\,MHz, $\mu=1$, DSUUU. \texttt{latency\_bound}: 8 interactive
5\,Mbps streams with a 12\,ms PDB sharing a saturated downlink with
80\,Mbps of bulk. Carrier: 40\,MHz, $\mu=1$, DDDSU (downlink-heavy).

The MCS staircase is a fitted $0.75\times$ Shannon curve capped at
7.5\,bits/RE ($\approx$5.8\,bps/Hz effective) at 31\,dB. A typical
factory UE at 22\,dB sits at 4.5\,bits/RE; a cell-edge UE at 14\,dB at
2.0. This is deliberately approximate---good enough for the
comparative claims here; a 3GPP TBS extract is future work.

The load sweep varies offered load around the shipped operating point.
Because capacity in a real deployment is fixed by spectrum, we report the
axis as \emph{load} relative to the shipped point, where higher is worse.
Uplink runs with an 8-slot ($\approx$4\,ms) BSR round-trip and an 8-slot
CQI report delay throughout.

\subsection{Metrics}
We report \textbf{contract-oriented} metrics, and this is a deliberate
methodological choice: mean delivery ratio hid every finding in this study.
A GBR flow's contract is its GFBR, counted as met at $\ge 95\%$; a Delay
flow's is $\ge 99\%$ on-time delivery within PDB. We report the
\emph{count of flows meeting their contract}, the \emph{minimum} delivery
across the GBR set, and p99 head-of-line delay. Mean and total throughput
appear as context, not as the verdict---a scheduler can post a healthy 86\%
mean while the missing 14\% is one starved safety-critical flow.

% ===========================================================================
\section{Results}\label{sec:results}
% ===========================================================================

\subsection{The value of QoS-awareness is a hump}\label{sec:hump}

Table~\ref{tab:sweep} sweeps offered load. At the shipped load GBR demand
far exceeds capacity and essentially no scheduler honours a contract; at one
third of it everyone has slack and all converge. In between, the two designs
differ in \emph{which} metric they favour rather than in which is better. At
$0.50\times$ load the two-tier design honours \textbf{10/10} contracts
against PF's 8/10; at $0.67\times$ PF meets 6/10 to our 1/10 while its
worst-served flow sits at 21\% against our 67\%. Section~\ref{sec:adopt}
takes that trade apart.

\begin{table}[t]
\caption{Load sweep on \texttt{factory\_robots}. Load is relative to the
shipped operating point; higher is worse. ``met'' counts GBR flows
delivering $\ge 95\%$ of GFBR. $-$mm pins the max-min stage off.}
\label{tab:sweep}
\centering
\small
\begin{tabular}{@{}llrrrr@{}}
\toprule
Load & Scheduler & Total & met & mean & min \\
\midrule
\multirow{3}{*}{$1.00\times$}
 & PF               & 69.3\,M & 1/10 & 51\% & 1\% \\
 & TwoTier          & 65.9\,M & 0/10 & 49\% & \textbf{44\%} \\
 & TwoTier $-$mm    & \textbf{73.1}\,M & 1/10 & 57\% & 0\% \\
\midrule
\multirow{3}{*}{$0.67\times$}
 & PF               & \textbf{96.5}\,M & \textbf{6/10} & 69\% & 21\% \\
 & TwoTier          & 93.2\,M & 1/10 & 71\% & \textbf{67\%} \\
 & TwoTier $-$mm    & 95.1\,M & 2/10 & 74\% & 47\% \\
\midrule
\multirow{2}{*}{$0.50\times$}
 & PF               & 116.8\,M & 8/10 & 80\% & 52\% \\
 & TwoTier          & \textbf{119.4}\,M & \textbf{10/10} & \textbf{83\%} & \textbf{81\%} \\
\midrule
\multirow{2}{*}{$0.33\times$}
 & PF               & 124.8\,M & 10/10 & 85\% & 83\% \\
 & TwoTier          & 125.5\,M & 10/10 & 85\% & 83\% \\
\bottomrule
\end{tabular}
\end{table}

The $0.67\times$ row is where the \emph{metric} becomes the discussion.
The two-tier design meets \emph{fewer} knife-edge contracts than PF (0/10
against 4/10) while delivering a strictly better distribution: minimum 60\%
against 4\%, mean 65\% against 55\%. A 95\% threshold rewards a scheduler
that abandons some flows so others clear the bar. For a factory where every
robot matters, ``all robots degraded gracefully'' beats ``four robots
perfect, six offline''---so the contract count is necessary but not
sufficient, and must be read alongside the minimum.

\textbf{Engineering implication.} Since capacity is fixed by spectrum, the
operator's lever is admission control: keep peak offered load in the
moderate band, because that is where the scheduler pays for itself. A cell
that systematically runs at ${\gtrsim}2.5\times$ its capacity has a
capacity-planning problem that no MAC fixes.

\subsection{Control-channel-limited: a capability PF lacks}\label{sec:pdcch}

Table~\ref{tab:sensor} reports \texttt{sensor\_dense}, where the DCI/CCE
budget binds before the data channel. The two-tier design meets
\textbf{30/30} deadlines with a 5\,ms p99 tail; PF meets 2/30 and pins at
the 15\,ms PDB. This is not a tuning difference. Each periodic flow
receives a standing allocation that costs \emph{zero PDCCH per slot} and
\emph{needs no BSR round-trip}; PF is throttled by both budgets at once and
has no mechanism to bypass either.

\begin{table}[t]
\caption{\texttt{sensor\_dense}: 30 periodic uplink sensors, 15\,ms PDB.
On-time requires $\ge 99\%$ delivery within PDB.}
\label{tab:sensor}
\centering
\small
\begin{tabular}{@{}lrrrr@{}}
\toprule
Scheduler & Total & on-time & min deliv. & p99 \\
\midrule
RoundRobin & 6.9\,M & 0/30  & 71\%  & 15.0\,ms \\
PF         & 9.2\,M & 2/30  & 85\%  & 15.0\,ms \\
\textbf{TwoTier}  & \textbf{9.6}\,M & \textbf{30/30} & \textbf{100\%} & \textbf{5.0}\,ms \\
\bottomrule
\end{tabular}
\end{table}

We regard this as the strongest practical result in the paper: for any
deployment with dense periodic small-payload traffic---sensors, PLCs, AGV
telemetry---configured-grant support is not an optimisation but a
prerequisite, and capacity planning that counts only PRBs will mis-size the
cell.

\subsection{Deadline-blindness is silent}\label{sec:latency}

Table~\ref{tab:latency} reports \texttt{latency\_bound}. The two-tier
design meets 8/8 deadlines; PF meets 5/8. The trade is explicit and
deliberate: bulk throughput falls from 24.6\,Mbps to 13.2\,Mbps to fund the
interactive set.

\begin{table}[t]
\caption{\texttt{latency\_bound}: 8 interactive streams (12\,ms PDB) versus
80\,Mbps of bulk on a saturated downlink.}
\label{tab:latency}
\centering
\small
\begin{tabular}{@{}lrrrr@{}}
\toprule
Scheduler & on-time & mean & p99 & bulk DL \\
\midrule
RoundRobin  & 3/8 & 84\%  & 12.0\,ms & 22.8\,M \\
PF          & 5/8 & 86\%  & 12.0\,ms & \textbf{24.6}\,M \\
\textbf{TwoTier} & \textbf{8/8} & \textbf{100\%} & \textbf{4.5}\,ms & 13.2\,M \\
\bottomrule
\end{tabular}
\end{table}

The dangerous property is that PF's failure is \emph{quiet}. Its
control-flow mean delivery is 86\%---a number that reads as healthy on a
dashboard---but the missing 14\% are aged-out packets, and in a
teleoperation or motion-control loop those are precisely the
safety-relevant ones. Mean delivery is not a safety metric.

\subsection{Cell-edge starvation and the max-min floor}\label{sec:celledge}

Table~\ref{tab:perflow} gives the per-flow picture at the shipped load and
makes the knapsack result of Section~\ref{sec:knapsack} concrete. Without
the max-min stage the two lowest-SE flows land at exactly $0\%$---not
degraded, abandoned---while high-SNR peers run at 71--80\%. With the stage
enabled no GBR flow falls below 53\%, against PF's worst at 1\%. Note the
aggregate cost: PF and RR both carry a \emph{higher} mean here, and the case
for the design at this load is entirely distributional.

\begin{table}[t]
\caption{Per-flow delivery as a fraction of GFBR, \texttt{factory\_robots}
at the shipped load, worst channel first. $-$mm pins the max-min stage off.}
\label{tab:perflow}
\centering
\small
\begin{tabular}{@{}lrrrrr@{}}
\toprule
Flow & SNR & RR & PF & TwoTier $-$mm & \textbf{TwoTier} \\
\midrule
ue7  & 14\,dB & 23\% & 28\% & \textbf{0\%} & 54\% \\
ue4  & 16\,dB & 51\% & 63\% & \textbf{0\%} & 53\% \\
ue9  & 17\,dB & 90\% & 83\% & 55\% & 55\% \\
ue2  & 18\,dB & 57\% & 68\% & 53\% & 54\% \\
ue6  & 19\,dB & 37\% & 46\% & 50\% & 55\% \\
ue3  & 20\,dB & 63\% & 74\% & 69\% & 54\% \\
ue10 & 20\,dB & 92\% & 1\%  & 69\% & 59\% \\
ue8  & 21\,dB & 100\% & 100\% & 69\% & 89\% \\
ue1  & 22\,dB & 77\% & 92\% & 71\% & 56\% \\
ue5  & 24\,dB & 56\% & 68\% & 80\% & 59\% \\
\midrule
mean &        & \textbf{65\%} & 62\% & 59\% & 59\% \\
\bottomrule
\end{tabular}
\end{table}

The cost is stated plainly: at the shipped load PF carries
\emph{more} total throughput than the two-tier design (69.3 vs
66.7\,Mbps). The max-min floor gives up 4\% of aggregate capacity to hold
the worst-served flow at 40\% of contract rather than 0\%. Whether that is
the right trade is a deployment question---a robot whose video is switched
off entirely is a failure in a way that a uniformly degraded fleet is
not---but it is a trade, and we do not present it as a free win.

Two further effects are visible. The mixed-flow UEs (ue8--ue10, carrying a
GBR flow \emph{and} a best-effort flow) collapse to 0--3\% under the
QoS-blind baselines, because their MAC multiplexer fills the transport
block by raw backlog and a continuously backlogged best-effort flow wins
every time. And the bill for the max-min floor is paid by the high-SNR
flows: ue1 drops from 91\% under PF to 56\%.

\subsection{Sensitivity to control-plane fidelity}\label{sec:sensitivity}

Because our BSR and CQI models are approximations, we sweep them. BSR delay
matters and is close to linear: over 0--16 slots PF's minimum delivery
falls 66\%$\to$30\% and its contract count breaks 8/10$\to$5/10, while the
two-tier design moves 80\%$\to$77\% and 10/10$\to$8/10 and loses about 1\%
of throughput against PF's 10\%. SPS-served flows are invariant to BSR by
construction. BSR \emph{loss} is near-inert on this workload (0--20\% loss
moves nothing) because continuously backlogged video keeps buffers
non-empty, so a lost report rarely changes the eligibility bit.

CQI staleness is likewise small here, and the reason is a property of the
channel rather than of the scheduler: our AR(1) model produces per-slot SNR
innovations well below the ${\sim}3$\,dB spacing of MCS thresholds, so even
a stale report usually picks the right MCS. A conservative semi-persistent
MCS for SPS---a mobility hedge---is neutral at 0\,dB margin and actively
harmful beyond 1\,dB, because larger reservations cross the SPS viability
floor. We report this as a \emph{negative} result for our deployment class:
the hedge is worth paying for under mobility, and a static factory is
exactly where it is not.

\subsection{Two negative results}\label{sec:negative}

We implemented and rejected two mechanisms, both aimed at cell-edge
starvation. An \textbf{adaptive dual-ascent penalty} escalates $p_i$ on
whichever flow is actually missing. It raises minimum delivery
(0\%$\to$34\%) but meets \textbf{0/10} contracts in both overloaded rows,
because dual ascent drives toward equal \emph{normalised shortfall} while a
GBR contract is a \emph{step function}---equalising shortfall parks every
flow just below its floor so none clears it. A
\textbf{spectral-efficiency tilt} $p_i \propto \mathrm{SE}_i^{k}$ with
$k<0$ rescues the cell edge only by relocating starvation to the next-worst
tier.

Section~\ref{sec:knapsack} explains why neither could have worked: both
reweight a linear penalty whose optimum is a vertex, so they select a
different victim rather than removing the abandonment. We think this is the
transferable lesson---\emph{in an infeasible allocation problem, a weight
chooses whom to sacrifice; only a constraint prevents sacrifice}---and it
applies to any scheduler enforcing rate contracts through penalty terms.

% ===========================================================================
\section{Three Modelling Shortcuts That Changed Our Conclusions}\label{sec:fidelity}
% ===========================================================================

Every result in Section~\ref{sec:results} is a claim about a scheduler
mediated by a simulator. Three times during this work a natural-looking
modelling shortcut produced a result that did not survive being corrected.
We report them with the corrected measurement beside each, because the
shortcuts are ones any scheduler evaluation might make, and because in one
case the artifact was a headline finding of an earlier draft of this paper.

\subsection{The scheduler filled the uplink transport block}

Our multiplexer was direction-agnostic: the same code split a UE's transport
block in both directions, ordering flows by
$(\texttt{priority}, -\tilde{Q})$ with $\tilde{Q}$ the \emph{gNB's} virtual
queue. In the uplink no gNB may do that. It grants a block; the UE fills it
by its own logical-channel prioritisation (TS 38.321
\S5.4.3.1)---two rounds in priority order, the first bounded by each
channel's prioritised-bit-rate (PBR) token bucket. The gNB configures those
PBRs but does not run the algorithm and never learns the split.

The shortcut handed our scheduler a per-slot lever on the uplink byte split
that does not exist, and it produced a result we believed:

\begin{center}\small
\begin{tabular}{@{}lrrrr@{}}
\toprule
mixed-flow UE & RR & PF & Gradient & TwoTier \\
\midrule
scheduler fills TB & 3\% & 3\% & 3\% & \textbf{53\%} \\
\textbf{UE fills TB} & \textbf{100\%} & \textbf{100\%} & 72\% & 89\% \\
\bottomrule
\end{tabular}
\end{center}

Under the shortcut the QoS-blind baselines collapse a GBR video flow to
${\sim}3\%$ of contract when a greedy best-effort flow shares its UE, while
our design holds it near 53\%. We described that as a clean demonstration of
QoS-aware multiplexing. It was not: with the UE filling its own block, the
baselines reach 100\% unaided. \textbf{The protection was the UE's PBR
configuration}---available to every scheduler, and nothing to do with ours.
Correcting it also raised PF's contract count at $0.50\times$ load from 5/10
to 8/10, which is most of the margin an earlier draft claimed.

\subsection{Flows shared a priority level}

With every flow at one priority, the multiplexer's priority sort is a stable
sort over equal keys and silently falls back to the order flows appear in the
scenario file. Results were reproducible but only accidentally: reordering a
configuration file would have changed them. Deriving each flow's priority
from its standardised 5QI (TS 23.501 Table 5.7.4-1) fixed it. Notably the
numbers \emph{did not change}---the stable sort had been producing the 5QI
ordering by luck---which is precisely why this class of error survives
testing.

\subsection{Tier-1 was handed the true offered load}

Our strategic tier read each flow's demand from the traffic generator's own
parameters: exact, noise-free, and static. No gNB has that; it must infer
demand from what it observes. Replacing the oracle with an estimator built
from arrivals produced a \emph{starvation lock-in}---a flow that lost out
appeared to be asking for less, so its cap fell, so it lost more. Six
contending flows went from $(3.9\ldots7.0)$\,Mbps to $(0.2\ldots11.8)$,
Jain's index $0.97 \to 0.47$.

The cause was subtle and worth naming: our estimator counted
delivered${}+{}$backlog, omitting bytes the flow \emph{discarded} on PDB
expiry. A starved flow's data ages out, so it stops appearing in the
arrival count while its true offered rate is unchanged. The flow never
stopped asking; we stopped counting.

The resolution was better than a fixed estimator. Tier-2's windowed ceiling
already clamps a flow's virtual queue to what actually arrived, so a quiet
flow is never scheduled however generous its target---and it does this from
\emph{observation} where the Tier-1 cap needed a \emph{prediction}. The two
were redundant, and the redundant copy was the one sitting where the
information is worst. Removing the demand cap reproduces oracle performance
exactly ($32.4$\,M, Jain $0.97$) with no demand estimate at all.

\subsection{What we take from this}

Each of the three gave the scheduler information or authority the 3GPP split
does not grant it, and each flattered the scheduler rather than the
baselines. That direction is not coincidence: a simulator is written from the
scheduler's point of view, so its conveniences accrue there. The check we now
apply is to ask, of every quantity the scheduler reads, \emph{which network
element learns this, and how}.

% ===========================================================================
\section{Adoption: What We Decided and Why}\label{sec:adopt}
% ===========================================================================

Two mechanisms win unconditionally: configured grants and the deadline-aware
tracker (Section~\ref{sec:results}). Both address capabilities PF lacks
entirely, and---the reason we trust them---both survived every correction in
Section~\ref{sec:fidelity} untouched, because neither scenario contains a
multi-flow uplink UE.

The GBR picture is mixed, and the summary ``PF wins in overload'' is wrong.
At $0.67\times$ load PF meets 6/10 contracts to our 1/10, but delivers a
worst-served flow at 26\% against our 82\%, and a lower mean (83\% against
89\%). Every robot above 82\% of its contracted rate is a better outcome than
six robots perfect and the worst at 26\%---unless partial delivery is
worthless, which for video and telemetry it is not and for a hard control
loop it is. \textbf{The metric, not the scheduler, is what differs}: a
threshold count rewards concentrating a shortfall, and this design spreads
it.

We tested whether a fallback could recover the difference. Sweeping the
max-min scale from 1 to 0 moves the contract count only from 1/10 to 2/10,
against PF's 6/10---so the gap is not the max-min floor but the fact that
Tier-1 allocates toward targets while PF's opportunism concentrates.
Recovering it would need contract \emph{selection}, which is admission
control.

A scheduler-level switch to PF is therefore the wrong instrument: it would
surrender configured grants and deadline awareness to improve one metric on
one flow class. The mechanisms are independent and only the max-min floor is
regime-dependent. It is also a continuous knob with a monotone effect, and
the strategic tier already computes $t^\star$---an \emph{exact} feasibility
test for the GBR set, not a heuristic---so degradation can be graceful and
needs no regime detection. Nor is detection needed for safety: at
$\mathrm{scale}=0$ the design leads PF on mean and worst-case delivery at
both loaded points while keeping both structural wins.

% ===========================================================================
\section{Threats to Validity}\label{sec:threats}

\textbf{Simulation, not silicon.} Results are comparative and
architecture-level: they establish which scheduler wins and roughly by how
much, not absolute throughput, microsecond latency, real-time feasibility or
UE interop. An OAI integration is the necessary next step, and the
\texttt{scheduler/} library is dependency-isolated for it.

\textbf{Absolute figures are soft.} A single run reads GBR delivery about 5
points below steady state with $\pm$2 point scatter; we verified over 60
consecutive windows that the scheduler-versus-scheduler \emph{gap} is
stable, so comparative claims hold. Spectral efficiency is a fitted
staircase rather than a 3GPP TBS extract.

\textbf{Scenario coverage.} Three synthetic scenarios isolate three
bottlenecks cleanly, but they are synthetic, single-cell, and have a fixed UE
population---no churn, no mobility, no inter-cell interference. Trace-driven
factory workloads would strengthen external validity considerably.

\textbf{Known unmodelled effects are conservative.} UL $k_2$
grant-to-transmission timing and HARQ retransmissions consuming PRBs are both
absent, and both would \emph{widen} the configured-grant advantage rather
than narrow it. Given Section~\ref{sec:fidelity}, we note that this
direction is asserted from the mechanism, not measured.

\textbf{A contract-dimensioning bound.} Even at three times capacity about
15\% of GBR video bytes expire on PDB. We verified this is
scheduler-independent---a lone video flow drops identically under RR, PF and
our design---because an I-frame burst is 3--4$\times$ the GFBR, making a
``GFBR plus tight PDB'' contract internally inconsistent. It bounds what any
scheduler can achieve on such a workload.

\section{Conclusion}

We set out to decide whether to replace the proportional-fair scheduler our
private-5G stack inherits from OpenAirInterface. The answer is yes, and the
reasoning is worth more than the verdict.

Two mechanisms justify the change on their own, and neither is a tuning
gain: configured grants, which meet 30/30 deadlines where PF meets 2/30 once
the control channel binds, and a deadline-aware per-slot tracker, without
which PF's misses are invisible in mean throughput. Both reproduce results
the field has reported separately, in a simulator sharing no code with the
work that first reported them. On guaranteed-bitrate contracts the picture is
mixed and the honest summary is that the \emph{metric} differs rather than
the outcome: a threshold count rewards concentrating a shortfall, this design
spreads it, and which is preferable depends on whether partial delivery has
value---for video and telemetry it does, for a hard control loop it does not.

The part we expect to be most useful to others is Section~\ref{sec:fidelity}.
Three natural modelling shortcuts produced results in our own evaluation that
did not survive correction, and in the worst case a headline finding
disappeared entirely: what we had described as our scheduler protecting
multi-flow uplink UEs turned out to be the UE's prioritised-bit-rate
configuration, visible only once the uplink transport block was filled by the
UE rather than by us. All three shortcuts flattered the scheduler over the
baselines, which we do not think is coincidence---a simulator is written from
the scheduler's point of view.

We make no novelty claim for the design itself; it composes established
parts, and we have described it because the composition and its failure modes
are instructive. What we can offer is a decision taken carefully, the
measurements behind it, the errors found on the way, and the code to check
any of it.

\section*{Artifact Availability}
% Double-blind: the repository is deliberately unnamed here. Restore the URL
% for camera-ready.
The simulator, the scheduler library, all scenario configurations and the
scripts named below will be released publicly on acceptance, and are
available to reviewers on request through the programme chairs. Every
quantitative claim in this paper is regenerated by a single command:
Tables~\ref{tab:sweep}, \ref{tab:sensor} and~\ref{tab:latency} by
\texttt{scheduler\_study.py}; Table~\ref{tab:perflow} by
\texttt{maxmin\_study.py}; the staircase, penalty sweep and utility-only
control of Section~\ref{sec:knapsack} by \texttt{knapsack\_diagnostic.py}; the sensitivity results of
Section~\ref{sec:sensitivity} by \texttt{bsr\_study.py} and
\texttt{cqi\_study.py}. A dated engineering log recording every finding,
regression and rejected design behind this paper is included.

\bibliographystyle{IEEEtran}
\bibliography{refs}

\end{document}
